%% ma: L12-B14-0026
%% lop: 12
%% chuong: 5
%% bai: 14
%% dang: 12.14.17
%% loai: TN4
%% muc: TH
%% dap_an: "B"
%% nguon: "CĐ2.5.pdf giáo viên gửi 27/10/2026 (M07, Lê Minh Khôi – Nguyễn Anh Khôi, Chương II Chủ đề 5 Tích có hướng của hai vectơ), Câu 10"
%% kien_thuc: "Ba vectơ đồng phẳng, tích hỗn tạp (chương trình cũ, không có trong SGK KNTT 12)"
%% trang_thai: cach-ly
%% ly_do: "Vượt chương trình: điều kiện đồng phẳng của ba vectơ (tích hỗn tạp) không có trong SGK KNTT 12. Gợi ý: bỏ, hoặc đổi thành câu bốn điểm đồng phẳng giải bằng phương trình mặt phẳng."
\begin{ex}
Trong không gian $Oxyz$, cho ba vectơ $\vec a=(m+2;3;2m)$, $\vec b=(2;-1;m)$ và $\vec c=(1;2;1)$ với $m\in\mathbb R$ là tham số. Gọi $S$ là tập hợp các giá trị của tham số $m$ để ba vectơ trên đồng phẳng. Số phần tử của tập hợp $S$ bằng
\choice{$0$}{\True $1$}{$2$}{$3$}
\loigiai{$\left[\vec b,\vec c\right]=(-1-2m;m-2;5)$, $\vec a\cdot\left[\vec b,\vec c\right]=-(m+2)(2m+1)+3(m-2)+10m=-2m^2+8m-8=0\Leftrightarrow m=2$. Có $1$ giá trị.}
\end{ex}
